\documentclass[10pt,twocolumn]{article}

% ---- Packages (all in Overleaf's default TeX Live) ----
\usepackage[utf8]{inputenc}
\usepackage[T1]{fontenc}
\usepackage[margin=0.85in]{geometry}
\usepackage{booktabs}
\usepackage{amsmath,amssymb}
\usepackage{graphicx}
\usepackage{xcolor}
\usepackage{authblk}
\usepackage{titlesec}
\usepackage[hidelinks]{hyperref}

\titleformat{\section}{\normalfont\large\bfseries}{\thesection}{0.5em}{}
\titleformat{\subsection}{\normalfont\normalsize\bfseries}{\thesubsection}{0.5em}{}
\setlength{\parskip}{2pt}

\title{\textbf{Phase-Conditioned Cascading Memory Routing\\ for Multi-Agent LLM Systems}}
\author[1]{[Your Name]}
\author[1]{[Mentor Name]}
\affil[1]{[Affiliation] \quad \texttt{[email]}}
\date{}

\begin{document}
\maketitle

\begin{abstract}
\noindent
Memory-augmented LLM agents increasingly maintain several specialized memory
stores---procedural rules, working state, cached plans, episodic traces,
semantic facts, and durable entity profiles. Yet most systems retrieve from
\emph{all} available stores on \emph{every} step, paying retrieval cost and
injecting irrelevant context regardless of what the current task needs. We
present the \textbf{Phase-Conditioned Cascading Memory Router (PCCR)}, a single
central manager that routes both reads and writes across six cognitive memory
types over a ten-phase agentic lifecycle, in a multi-agent
orchestrator/sub-agent setting. PCCR (i)~makes the routing policy a function of
the \emph{lifecycle phase} being executed, (ii)~on a cache miss gates each
optional store by an explicit cost-effectiveness ratio
$\rho = \text{utility}/\text{cost}$ against a tunable threshold, and
(iii)~logs every decision as a first-class record that a closed loop uses to
adapt utilities from task outcomes. On the OfficeBench office-automation
benchmark with a real backbone (gpt-oss-120b), PCCR matches or exceeds the
accuracy of an indiscriminate \emph{retrieve-everything} policy while issuing
\textbf{41--87\% fewer store consultations} and injecting up to \textbf{89\%
fewer retrieved tokens} at equal accuracy, in both two-agent (calendar+email)
and three-agent (calendar+email+document) configurations. A per-task-pattern
analysis shows the savings arise from interpretable, store-specific
decisions---e.g.\ episodic search is pruned exactly on pure-lookup tasks where
it cannot help. We position PCCR against recent agent-memory routers and ground
its score in classical value-of-information theory.
\end{abstract}

\section{Introduction}
LLM agents lose all state when a task ends unless augmented with external
memory. Cognitive-architecture accounts of agents (CoALA~\cite{coala}) and
recent surveys~\cite{survey1,survey2} converge on a small set of functionally
distinct memory types---\emph{procedural} (skills/rules), \emph{working}
(volatile task state), \emph{short-term} (a session cache), \emph{episodic}
(specific past experiences), \emph{semantic} (distilled facts), and
\emph{entity} (durable facts about users/things). A practical agent benefits
from all of them.

This raises a routing problem that is largely unaddressed: \emph{given many
memory stores of differing cost, which should the agent consult for the current
task, and which should it write to?} The prevailing answer in deployed systems
is \emph{all of them, every time}. This is wasteful: a similarity search over a
large episodic store is orders of magnitude costlier than a key--value lookup of
a user profile, and for many tasks the expensive store contributes
nothing---a factual lookup (``when is my meeting with Alice?'') is answered from
stored facts, not from a search through past task traces.

Recent work has begun to treat retrieval as a decision rather than a reflex.
Cost-sensitive store routing~\cite{pocket} frames store selection as an
accuracy-vs-cost trade-off and shows (via an oracle) that selective retrieval
dominates uniform retrieval; learned budget-tier routers~\cite{budgetmem} and
embedding-based write routers~\cite{memrouter} optimize specific edges of the
problem. However, each existing system optimizes \emph{one} edge---write
admission, \emph{or} read-side store selection, \emph{or} memory-structure
choice, \emph{or} a context budget---under a \emph{single, global} policy, and
almost always in a \emph{single-agent conversational} loop. None route across a
full six-type cognitive taxonomy, with one central manager governing both reads
and writes, conditioned on \emph{which lifecycle phase is executing} and
\emph{what kind of task it is}, in a \emph{multi-agent} orchestrator/sub-agent
system.

\paragraph{Contributions.} We close that gap with PCCR:
\begin{enumerate}\itemsep1pt
\item \textbf{Phase-conditioned routing.} The policy is not global; each of ten
lifecycle phases carries its own rule, because ``what to consult'' means
something structurally different at retrieval time than at storage or
consolidation time (\S\ref{sec:phase}).
\item \textbf{A cost-effectiveness cascade.} Retrieval first attempts a free
short-circuit (a short-term cache hit skips all optional stores); on a miss it
consults an optional store \emph{iff} $\rho = U(\text{pattern},s)/C(s)$ clears a
single tunable threshold $\theta$ (\S\ref{sec:gate}).
\item \textbf{Decision logging + a closed loop.} Every routing call emits a
structured \texttt{RoutingDecision}; an outcome-driven update nudges utilities
up/down from success/failure (\S\ref{sec:loop}).
\item \textbf{Empirical evidence} on OfficeBench with a real backbone, in two-
and three-agent settings, that PCCR preserves accuracy while substantially
reducing retrieval cost (\S\ref{sec:results}).
\end{enumerate}

\section{Related Work}
\paragraph{Taxonomies and lifecycle.} CoALA~\cite{coala} maps human memory types
onto LLM agents; surveys~\cite{survey1,survey2} formalize agent memory as a
write--manage--read loop. These define \emph{what} the memory types are but not
\emph{how to route} among them.

\paragraph{Cost-sensitive / learned routing.} ``Did You Check the Right
Pocket?''~\cite{pocket} formalizes store selection as a cost-sensitive decision
and shows an oracle advantage, but gives no implementable policy and considers
retrieval in isolation. BudgetMem~\cite{budgetmem} learns per-module budget
tiers under an RL cost objective; U-Mem~\cite{umem} scores retrieval by
similarity plus a learned utility; AgeMem~\cite{agemem} learns tool-based
long/short-term management. These learn a policy over a \emph{flat} memory space
in single-agent settings, without phase conditioning.

\paragraph{Write-side, structure, multi-agent.} MemRouter~\cite{memrouter} is a
write-admission router; FluxMem~\cite{fluxmem} selects a memory \emph{structure}.
LEGOMem~\cite{legomem} splits \emph{procedural} memory into orchestrator vs.\
sub-agent modules (informing our PM split) but has no central cost-gated router;
MIRIX~\cite{mirix} defines six memory types managed by \emph{separate} per-type
agents, with no single central cost-based manager; RCR-Router~\cite{rcr} routes a
context \emph{budget} by role within one phase.

\paragraph{Consolidation, production, grounding.} A-MEM~\cite{amem} and
self-consolidation~\cite{selfcons} inform our consolidation phase;
MemGPT~\cite{memgpt}, Mem0~\cite{mem0}, MemMachine~\cite{memmachine}, and
MemOS~\cite{memos} apply uniform policies without phase/pattern-conditioned cost
gating. The cost/utility score descends from value-of-information
theory~\cite{voi} and metareasoning~\cite{meta}. To our knowledge no prior system
combines a six-type taxonomy, a single central read+write router, phase +
task-pattern conditioning, and a multi-agent setting.

\section{The Agentic Memory Lifecycle}
We model a task as flowing through \textbf{ten phases}: bootstrap, ingestion,
retrieval, orchestrator inference, delegation, agent inference, agent output,
observation, storage, and consolidation. Six memory types carry data across
them (Table~\ref{tab:types}). Each task is classified at ingestion into a
\textbf{pattern} (lookup, single-action, coordination, recurring,
document-create, document-query, \ldots) by a lightweight classifier; this label
conditions all downstream routing.

\begin{table}[t]
\centering\small
\begin{tabular}{@{}lll@{}}
\toprule
Type & Role & Access cost \\
\midrule
PM (proc.)   & static prompts/rules        & cheap (dict) \\
WM (work.)   & per-task scratchpad          & cheap \\
STM (short)  & cache of plan bundles        & cheap \\
EM (episodic)& past task traces (FAISS)     & \textbf{expensive} \\
SM (semantic)& distilled facts (FAISS)      & \textbf{expensive} \\
ENT (entity) & per-user profiles            & cheap \\
\bottomrule
\end{tabular}
\caption{The six memory types and their relative access cost.}
\label{tab:types}
\end{table}

\section{Method: PCCR}
A central \texttt{MemoryManager} executes the lifecycle; a \texttt{MemoryRouter}
\emph{decides} what to touch at each phase, and the manager performs the
reads/writes. This decision/execution separation lets the policy be inspected,
tuned, and replaced without rewriting the lifecycle plumbing.

\subsection{Phase conditioning}\label{sec:phase}
Each phase declares a routing \emph{mode}: bootstrap $=$ classify-by-kind
(rules$\rightarrow$PM, facts$\rightarrow$ENT, vectors$\rightarrow$SM); ingestion
and observation $=$ fixed-sink ($\rightarrow$WM); retrieval $=$ the cascading
cost-gate (below); orchestrator inference $=$ consume-and-log (nothing persists);
storage $=$ outcome-conditional (success $\rightarrow$ write STM+ENT; failure
$\rightarrow$ write nothing); consolidation $=$ batch-distribute
($\rightarrow$EM/SM/PM/STM). The same request shape is routed differently
depending on the phase.

\subsection{The cascading cost-effectiveness gate}\label{sec:gate}
PM and WM are read unconditionally. STM is checked first; a confident cache hit
\textbf{short-circuits} all optional stores. On a miss, for each optional store
$s \in \{\text{EM},\text{SM},\text{ENT}\}$:
\[
\rho(\text{pattern},s) = \frac{U(\text{pattern},s)}{C(s)},\quad
\text{consult } s \iff \rho \geq \theta,
\]
where $C(s)$ is a relative access cost ($C(\text{ENT}){=}0.15$ for a keyed
lookup; $C(\text{EM}){=}C(\text{SM}){=}0.55$ for a vector search +
deserialization) and $U(\text{pattern},s)\in[0,1]$ is the expected utility of $s$
for that pattern. $\theta$ is a single dial trading cost against quality. $U$
encodes hypotheses such as ``pure-lookup tasks gain little from episodic search
but much from entity facts,'' which the experiments test directly.

\subsection{Outcome-driven adaptation and audit}\label{sec:loop}
Every routing call returns and records a \texttt{RoutingDecision}: candidates,
consulted set, skipped set with per-store reasons, cache hit/miss, and latency. A
closed loop reads this log and, after each task, nudges $U(\text{pattern},s)$
\textbf{up} if consulting $s$ coincided with success and \textbf{down} if with
failure (clamped to $[0,1]$). Because the policy is indexed by
$(\text{phase}\times\text{pattern}\times\text{store})$, the learned quantity is
higher-dimensional and more interpretable than a flat global policy.

\subsection{Priors, not assumptions}
$C(\cdot)$ and $U(\cdot,\cdot)$ start as hand-set priors. They are not assumed
final: the closed loop refines $U$ from outcomes, $C$ can be replaced by measured
latencies from the decision log, and---as \S\ref{sec:results} shows---the
held-out comparison itself empirically tests the priors (pruning a low-utility
store without losing accuracy confirms its low utility).

\section{Experimental Setup}
\paragraph{Benchmark.} OfficeBench~\cite{officebench}, a multi-application
office-automation benchmark; a task is correct only if \emph{all} of its
evaluation predicates pass. We use the subset executable by our agents: 34
calendar/email tasks (two-agent) and 92 calendar/email/document tasks
(three-agent).

\paragraph{Agents and backbone.} An orchestrator delegates to specialist
sub-agents---calendar, email, and (three-agent) a document agent. All roles run
on \textbf{gpt-oss-120b} (via Cerebras), temperature 0. EM/SM are FAISS
\texttt{IndexFlatIP} over 384-d \texttt{all-MiniLM-L6-v2} embeddings.

\paragraph{Protocol.} Held-out, disjoint train/test split (two-agent
$\sim$23/11; three-agent 60/32). Train (memory off) $\rightarrow$ consolidate
once $\rightarrow$ test with memory. The three-agent comparison uses
\textbf{frozen test memory} (no writes during test) so the only variable across
methods is the routing policy; training is run once and shared via
snapshot/restore.

\paragraph{Methods.} All share identical agents, prompts, and consolidated
memory, differing only in retrieval policy: \textbf{Retrieve-All} (consult every
optional store), \textbf{Boolean} (a hand-set per-pattern on/off table), and
\textbf{PCCR} (the $\rho$-gate, swept over $\theta\in\{1.0,1.2,1.4\}$).

\paragraph{Metrics \& RQs.} Task accuracy; optional-store consultations; FAISS
queries; retrieved tokens injected; utility drift. \textbf{RQ1}: cost at equal
quality? \textbf{RQ2}: does the consult set vary by pattern as predicted?
\textbf{RQ3}: is $\theta$ a monotone cost/quality dial? \textbf{RQ4}: does the
closed loop adapt?

\section{Results}\label{sec:results}

\subsection{RQ1 --- Cost at equal quality}
All methods reach identical accuracy; PCCR attains it at the lowest cost
(Tables~\ref{tab:cmp2}--\ref{tab:cmp3}).

\begin{table}[t]
\centering\small
\begin{tabular}{@{}lccc@{}}
\toprule
Method & Acc & Consults & FAISS \\
\midrule
Retrieve-All & 5/11 & 20 & 27 \\
Boolean      & 5/11 & 14 & 27 \\
\textbf{PCCR}& 5/11 & \textbf{12} & \textbf{25} \\
\bottomrule
\end{tabular}
\caption{Two-agent held-out comparison ($\theta{=}1.0$, 11 tasks).}
\label{tab:cmp2}
\end{table}

\begin{table}[t]
\centering\small
\begin{tabular}{@{}lccc@{}}
\toprule
Method & Acc & Consults & FAISS \\
\midrule
Retrieve-All & 10/32 & 46 & 72 \\
Boolean      & 10/32 & 39 & 72 \\
\textbf{PCCR}& \textbf{11/32} & \textbf{27} & \textbf{62} \\
\bottomrule
\end{tabular}
\caption{Three-agent held-out comparison ($\theta{=}1.0$, 32 tasks). PCCR
slightly exceeds the baselines' accuracy while cutting consultations
$\sim$41\%.}
\label{tab:cmp3}
\end{table}

\subsection{RQ3 --- Threshold as a cost dial; token savings}
The three-agent frontier (Table~\ref{tab:frontier}; all points share identical
consolidated memory) has a sharp knee: $\theta{=}1.0{\rightarrow}1.2$ cuts
consultations $27{\rightarrow}6$ ($-78\%$) and injected retrieved tokens
$837{\rightarrow}96$ ($-89\%$) with \emph{no} accuracy loss; accuracy only
degrades at $\theta{=}1.4$. \textbf{Headline: PCCR at $\theta{=}1.2$ matches or
beats every baseline's accuracy while issuing $\sim$87\% fewer consultations than
Retrieve-All (6 vs.\ 46) and injecting $\sim$89\% fewer retrieved tokens.} Token
counts are an orchestrator-context lower bound ($\sim$4 chars/token); realized
savings are larger.

\begin{table}[t]
\centering\small
\begin{tabular}{@{}ccccc@{}}
\toprule
$\theta$ & Acc & Consults & FAISS & Inj.\ tokens \\
\midrule
1.0 & 11/32 & 27 & 60 & 837 \\
\textbf{1.2} & \textbf{11/32} & \textbf{6} & 38 & \textbf{96 ($-89\%$)} \\
1.4 & 8/32 & 6 & 37 & 96 \\
\bottomrule
\end{tabular}
\caption{Three-agent cost/quality frontier (self-consistent).}
\label{tab:frontier}
\end{table}

\subsection{RQ2 --- Pattern-conditioned routing}
Per-pattern consult rates (Table~\ref{tab:pattern}) match the $\rho$-model
exactly: episodic is pruned on pure lookups (\texttt{email\_query},
$\rho_{\text{EM}}{=}0.87{<}1$), entity is pruned where a profile is irrelevant
(\texttt{email\_send}), both are consulted for coordination, and a pure document
read (\texttt{word\_query}) consults \emph{neither} optional store. The
aggregate savings are thus the sum of interpretable, store-specific decisions.

\begin{table}[t]
\centering\small
\begin{tabular}{@{}lccc@{}}
\toprule
Pattern & $n$ & EM & ENT \\
\midrule
\texttt{email\_query} (lookup)        & 7  & 0/7  & 7/7 \\
\texttt{email\_send}                  & 10 & 10/10& 0/10 \\
\texttt{single\_cal\_create}          & 8  & 8/8  & 0/8 \\
\texttt{multi\_cal\_find\_and\_create}& 1  & 1/1  & 1/1 \\
\texttt{word\_create}                 & 3  & 3/3  & 0/3 \\
\texttt{word\_query}                  & 3  & 0/3  & 0/3 \\
\bottomrule
\end{tabular}
\caption{Per-pattern EM/ENT consult rates (PCCR, three-agent, $\theta{=}1.0$).}
\label{tab:pattern}
\end{table}

\subsection{RQ4 --- Closed-loop adaptation}
In the online two-agent run, outcome-driven updates moved utilities
\textbf{bidirectionally}---upward for (pattern,\,store) pairs co-occurring with
success (e.g.\ \texttt{multi.episodic} $0.78{\rightarrow}0.83$) and downward for
those co-occurring with failure (e.g.\ \texttt{email\_send.episodic}
$0.62{\rightarrow}0.57$)---demonstrating online policy adaptation with no
backbone updates. (The three-agent run froze memory to isolate routing, so RQ4
is reported from the two-agent run.)

\section{Discussion and Limitations}
The pattern-level analysis (\S\,RQ2) is the strongest evidence: the mechanism
behaves as designed rather than succeeding by chance. \textbf{Scale:} held-out
sets are small (11/32 tasks); results are a directional proof-of-concept, and
some splits are pattern-skewed---scaling and balancing is the priority for a
camera-ready result. \textbf{Cost currency:} in our local harness FAISS is
milliseconds and LLM inference dominates wall-clock, so we report retrieval cost
(consultations, FAISS queries, injected tokens), the deployment-relevant cost
when stores are large, remote, or LLM-backed; token counts are approximate and a
lower bound. \textbf{Priors:} $C$ and $U$ begin hand-set (the loop refines $U$; $C$
is replaceable by logged latencies); a fully learned cost model and a
counterfactual (with/without-store) measurement of utility remain future work.

\section{Conclusion}
PCCR conditions retrieval on lifecycle phase and task pattern, gates each store
by an explicit cost-effectiveness ratio, and adapts from logged outcomes. On
OfficeBench with a real backbone, it matches or beats an indiscriminate
retrieve-everything policy while cutting store consultations by 41--87\% and
injected tokens by up to 89\% at equal accuracy, with savings attributable to
interpretable per-pattern decisions. ``Which memory to consult'' should be a
tunable, phase-aware, cost-sensitive decision rather than a reflex.

\begin{thebibliography}{99}\small
\bibitem{coala} T.~Sumers et al. Cognitive Architectures for Language Agents (CoALA). TMLR, 2024.
\bibitem{survey1} Memory for Autonomous LLM Agents: Mechanisms, Evaluation, and Emerging Frontiers. arXiv:2603.07670.
\bibitem{survey2} Anatomy of Agentic Memory: Taxonomy and Empirical Analysis. arXiv:2602.19320.
\bibitem{pocket} M.~Gaikwad. Did You Check the Right Pocket? Cost-Sensitive Store Routing for Memory-Augmented Agents. ICLR 2026 Workshop. arXiv:2603.15658.
\bibitem{budgetmem} H.~Zhang et al. Learning Query-Aware Budget-Tier Routing for Runtime Agent Memory (BudgetMem). arXiv:2602.06025.
\bibitem{memrouter} MemRouter: Memory-as-Embedding Routing for Long-Term Conversational Agents. arXiv:2605.00356.
\bibitem{umem} Towards Autonomous Memory Agents (U-Mem). arXiv:2602.22406.
\bibitem{agemem} Agentic Memory: Learning Unified Long-Term and Short-Term Memory Management for LLM Agents. arXiv:2601.01885.
\bibitem{fluxmem} Choosing How to Remember: Adaptive Memory Structures for LLM Agents (FluxMem). arXiv:2602.14038.
\bibitem{legomem} LEGOMem: Modular Procedural Memory for Multi-agent LLM Systems. arXiv:2510.04851.
\bibitem{mirix} Y.~Wang, X.~Chen. MIRIX: Multi-Agent Memory System for LLM-Based Agents. arXiv:2507.07957.
\bibitem{rcr} RCR-Router: Efficient Role-Aware Context Routing for Multi-Agent LLM Systems. arXiv:2508.04903.
\bibitem{amem} A-MEM: Agentic Memory for LLM Agents. NeurIPS, 2025. arXiv:2502.12110.
\bibitem{selfcons} Self-Consolidation for Self-Evolving Agents. arXiv:2602.01966.
\bibitem{memgpt} C.~Packer et al. MemGPT: Towards LLMs as Operating Systems. 2024.
\bibitem{mem0} Mem0: Building Production-Ready AI Agents with Scalable Long-Term Memory. 2025.
\bibitem{memmachine} MemMachine: A Ground-Truth-Preserving Memory System. arXiv:2604.04853.
\bibitem{memos} MemOS: An Operating System for Memory-Augmented Generation. EMNLP 2025. arXiv:2505.22101.
\bibitem{voi} R.~Howard. Information Value Theory. IEEE Trans. SSC, 1966.
\bibitem{meta} S.~Russell, E.~Wefald. Principles of Metareasoning. Artificial Intelligence, 1991.
\bibitem{officebench} T.~Wang et al. OfficeBench: Benchmarking Language Agents across Multiple Applications for Office Automation. arXiv:2407.19056.
\end{thebibliography}

\end{document}
